\documentclass[../../main.tex]{subfiles}% 注意这里的文件路径不能用 ./main.tex ,否则用latexmk编译子文件会报错
\graphicspath{{\subfix{./image/}}{\subfix{../../image/}}} % 指定图片引用路径,后续可以直接使用图片文件名
% 如果图片存放在上级目录,则使用 ../../image/ 作为备用路径

% 例如:
% \begin{figure}[H]
% \centering
% \includegraphics[width=0.3\textwidth]{图.png}
% \caption{}
% \label{figure:图}
% \end{figure}
% 注意:上述\label{}一定要放在\caption{}之后,否则引用图片序号会只会显示??.

\begin{document}

\section{常见问题}

\begin{example}
设的二次曲面方程为 $F(\bm{x}) = \bm{x}^T A \bm{x} + \bm{b}^T \bm{x} + c = 0,$ 其中 $\bm{x} = (x_1, x_2, \cdots, x_n)^T,$ $A$ 为 $n$ 阶实对称矩阵, $\bm{b}$ 为 $n$ 维列向量, $c$ 为常数. 其化为标准型的完整步骤如下：
\begin{enumerate}[(1)]
\item 正交变换消去交叉项：存在正交矩阵 $Q,$ 使得 $Q^T A Q = \Lambda = \mathrm{diag}(\lambda_1, \lambda_2, \cdots, \lambda_n).$ 令 $\bm{x} = Q \bm{y},$ 则方程化为见 \eqref{eq:quadratic_form_step_1}
\begin{align}
\bm{y}^T \Lambda \bm{y} + \bm{b}'^T \bm{y} + c = 0, \label{eq:quadratic_form_step_1}
\end{align}
其中 $\bm{b}' = Q^T \bm{b}.$
\item 平移变换消去一次项：若所有的特征值 $\lambda_i \neq 0,$ 令 $\bm{y} = \bm{z} + \bm{h},$ 其中 $\bm{h}$ 满足 $2 \Lambda \bm{h} + \bm{b}' = 0,$ 即 $\bm{h} = -\frac{1}{2} \Lambda^{-1} \bm{b}'.$ 此时常数项变为 $c' = c - \frac{1}{4} \bm{b}'^T \Lambda^{-1} \bm{b}'.$ 方程化为见 \eqref{eq:quadratic_form_step_2}
\begin{align}
\bm{z}^T \Lambda \bm{z} + c' = 0, \label{eq:quadratic_form_step_2}
\end{align}
\item 特例：若存在 $\lambda_i = 0$ (设前 $r$ 个非零, 后 $n-r$ 个为零), 则无法通过完全平移消去对应的一次项. 此时可将非零特征值的部分配方, 而零特征值的部分保留一次项. 方程化为见 \eqref{eq:quadratic_form_step_3}
\begin{align}
\sum_{i=1}^{r} \lambda_i z_i^2 + \sum_{j=r+1}^{n} d_j z_j + c' = 0, \label{eq:quadratic_form_step_3}
\end{align}
其中若存在 $d_j \neq 0,$ 则对应的曲面为抛物面或柱面.
\end{enumerate}
根据特征值 $\lambda_i$ 的符号以及常数项 $c'$ 的正负, 可分类判定曲面类型：若所有 $\lambda_i$ 同号且 $c' \neq 0$ 则为椭球面；若存在正负特征值且 $c' \neq 0$ 则为单叶或双叶双曲面；若 $c' = 0$ 则为二次锥面；若有 $\lambda_i = 0$ 且对应一次项不为零, 则为抛物面；若对应一次项也为零, 则为柱面.
\end{example}

\begin{example}
已知椭球面
$$
\Sigma_0:\frac{x^2}{a^2}+\frac{y^2}{b^2}+\frac{z^2}{c^2}=1,a>b,$$
的外切柱面$\Sigma_\varepsilon$($\varepsilon=1$或$-1$)平行于已知直线
$$
l_\varepsilon:\frac{x-2}{0}=\frac{y-1}{\varepsilon\sqrt{a^2-b^2}}=\frac{z-3}{c}.
$$
试求与$\Sigma_\varepsilon$交于一个圆周的平面的所有可能的法方向.
\end{example}
\begin{remark}
本题中的外切柱面指的是每一条直母线均与已知椭球面相切的柱面.
\end{remark}
\begin{solution}
设$M_1(x_1,y_1,z_1)$为$\Sigma_0$与$\Sigma_\varepsilon$的切点，则由条件可设$\Sigma_\varepsilon$上过$M_1$点的直母线为
\begin{align}
l:\begin{cases}
x=x_1\\
y=y_1+\varepsilon \sqrt{a^2-b^2}t\\
z=z_1+ct
\end{cases}, \quad t\in \mathbb{R}. \label{eq::::-=--2389tj34g34g3432r221.4}
\end{align}
将上式与$\Sigma_0$的方程联立得
\begin{align}
\frac{x_1^2}{a^2}+\frac{(y_1+\varepsilon \sqrt{a^2-b^2}t)^2}{b^2}+\frac{(z_1+ct)^2}{c^2}=1. \label{eq::::-=--2389tj34g34g3432r221.1}
\end{align}
又因为$M_1\in \Sigma_0$，所以
\begin{align}
\frac{x_1^2}{a^2}+\frac{y_1^2}{b^2}+\frac{z_1^2}{c^2}=1. \label{eq::::-=--2389tj34g34g3432r221.2}
\end{align}
再联立\eqref{eq::::-=--2389tj34g34g3432r221.1}\eqref{eq::::-=--2389tj34g34g3432r221.2}式可得
\begin{align*}
\frac{a^2}{b^2}t^2+\left( \frac{2\varepsilon \sqrt{a^2-b^2}}{b^2}y_1+\frac{2}{c}z_1 \right) t=0.
\end{align*}
由$l$与$\Sigma_0$相切可知上述方程只有$t=0$这一个解，故
\begin{align}
\frac{2\varepsilon \sqrt{a^2-b^2}}{b^2}y_1+\frac{2}{c}z_1=0. \label{eq::::-=--2389tj34g34g3432r221.3}
\end{align}
再将$\Sigma_0$的方程与\eqref{eq::::-=--2389tj34g34g3432r221.3}\eqref{eq::::-=--2389tj34g34g3432r221.4}式联立消去$x_1,y_1,z_1,t$得
\begin{align*}
\Sigma_\varepsilon:\frac{x^2}{a^2}+\frac{(cy-\varepsilon \sqrt{a^2-b^2}z)^2}{a^2c^2}=1.
\end{align*}
再设与$\Sigma_\varepsilon$交于一个圆周的平面为
\begin{align*}
\pi:Ax+By+Cz=0.
\end{align*}
记$K=\sqrt{A^2+B^2+C^2},M=\sqrt{A^2+B^2}$，现在做如下正交变换
\begin{align*}
\left[ \begin{array}{c}
x'\\
y'\\
z'\\
\end{array} \right] =\left[ \begin{matrix}
\frac{AC}{MK}&		\frac{BC}{MK}&		\frac{-A^2-B^2}{MK}\\
-\frac{B}{M}&		\frac{A}{M}&		0\\
\frac{A}{K}&		\frac{B}{K}&		\frac{C}{K}\\
\end{matrix} \right] \left[ \begin{array}{c}
x\\
y\\
z\\
\end{array} \right] 
\end{align*}
将$Oxyz$坐标系变为$Ox'y'z'$坐标系。在新坐标系$Ox'y'z'$下
\begin{align*}
\pi:z'=0.
\end{align*}
\begin{align*}
\Sigma _{\varepsilon}:\frac{\left( \frac{AC}{MK}x' -\frac{B}{M}y' +\frac{A}{K}z' \right) ^2}{a^2}+\frac{1}{a^2c^2}\left[ c\left( \frac{BC}{MK}x' +\frac{A}{M}y' +\frac{B}{K}z' \right) -\varepsilon \sqrt{a^2-b^2}\left( \frac{-A^2-B^2}{MK}x' +\frac{C}{K}z' \right) \right] ^2=1.
\end{align*}
故
\begin{align*}
\pi \cap \Sigma _{\varepsilon}:\frac{\left( \frac{AC}{MK}x' -\frac{B}{M}y' \right) ^2}{a^2}+\frac{1}{a^2c^2}\left[ c\left( \frac{BC}{MK}x' +\frac{A}{M}y' \right) -\varepsilon \sqrt{a^2-b^2}\left( \frac{-A^2-B^2}{MK}x' \right) \right] ^2=1.
\end{align*}
注意到交线$\pi \cap \Sigma_\varepsilon$为圆周的充要条件是上式中的$x'^2,y'^2$的系数相等且不为$0$，$x'y'$的系数为$0$。此即
\begin{align*}
\begin{cases}
\frac{2\varepsilon cA\sqrt{A^2+B^2}\sqrt{a^2-b^2}}{K}=0\\
c^2C^2+\left( a^2-b^2 \right) B^2+2\varepsilon c\sqrt{a^2-b^2}BC=c^2\left( B^2+C^2 \right) \ne 0\\
\end{cases},
\end{align*}
上式也等价于
\begin{align*}
\begin{cases}
A\sqrt{A^2+B^2}=0\\
BC=\frac{b^2+c^2-a^2}{2\varepsilon c\sqrt{a^2-b^2}}B^2\\
\end{cases},
\end{align*}
从而要么$A=B=0$,要么$A=0,B\ne 0,C=\frac{b^2+c^2-a^2}{2\varepsilon c\sqrt{a^2-b^2}}B.$
故$\pi$的法向量只可能为
\begin{align*}
(0,0,1)\text{或}(0,2\varepsilon c\sqrt{a^2-b^2},b^2+c^2-a^2).
\end{align*}

\end{solution}

\begin{example}
求二次曲面$C: x^2+y^2+z^2-4x-2y-4z+5=0$上的点到平面$\pi:4y-3z-7=0$的距离的最大值。
\end{example}
\begin{solution}
对$x^2+y^2+z^2-4x-2y-4z+5=0$配方得$(x-2)^2+(y-1)^2+(z-2)^2=4$。于是二次曲面$C$是以$(2,1,2)$为球心，$2$为半径的球面。球心到平面的距离为
\begin{align*}
d=\frac{|4\times 1-3\times 2-7|}{\sqrt{4^2+3^2}}=\frac{9}{5}.
\end{align*}
则二次曲面$C$上的点到平面$\pi:4y-3z-7=0$的距离的最大值为$d_{\mathrm{max}}=r+d=\frac{19}{5}$。
\end{solution}

\begin{example}
过点$P(0,0,2)$作与球面$x^2+y^2+z^2=1$相交的直线，记点$Q$为其交点连线的中点。证明：存在一定点，使得所有满足上述条件的点$Q$到该定点的距离为定值。
\end{example}
\begin{note}
解法二的思路是:先画图利用几何关系，猜出$(0,0,1)$是$Q$点轨迹的球心,再计算验证.实际上,可以通过选取两个不同交点的中垂线看出球心就是$(0,0,1)$。
\end{note}
\begin{solution}

{\color{blue}解法一:}
过$P$点的直线$L$与球面相交的交点分别为$A,B$，$Q(x,y,z)$，则$\overrightarrow{OQ}=(x,y,z)$，$\overrightarrow{PQ}=(x,y,z-2)$。由球面的性质$\overrightarrow{OQ}$与$AB$垂直，于是
\begin{align*}
\overrightarrow{OQ}\cdot \overrightarrow{PQ}=(x,y,z)\cdot(x,y,z-2)=x^2+y^2+(z-1)^2=1.
\end{align*}
即$Q$是以$(0,0,1)$为球心，$1$为半径的球面.再由直线$L$的任意性知，存在定点$(0,0,1)$，使点$Q$到该定点的距离为定值$1$。

{\color{blue}解法二:}
设过$P$点与球面相交的直线$L$方程为
\begin{align*}
L:\begin{cases}
x=at\\
y=bt\\
z=2+ct
\end{cases}\quad (t\in\mathbb{R}),
\end{align*}
其中$a,b,c\in\mathbb{R}$。
将上式与球面方程联立得
\begin{align*}
(a^2+b^2+c^2)t^2+4ct+3=0.
\end{align*}
由$L$与球面相交知上式必有解，记为$t_1,t_2$。设$Q$点坐标为$(at_Q,bt_Q,2+ct_Q)$，则
\begin{align*}
t_Q=\frac{t_1+t_2}{2}=-\frac{2c}{a^2+b^2+c^2},
\end{align*}
\[
Q\biggl(-\frac{2ac}{a^2+b^2+c^2},-\frac{2bc}{a^2+b^2+c^2},2-\frac{2c^2}{a^2+b^2+c^2}\biggr).
\]
从而$Q$到$M(0,0,1)$的距离平方为
\begin{align*}
|\overrightarrow{QM}|^2=\frac{4a^2c^2+4b^2c^2}{(a^2+b^2+c^2)^2}+\biggl(1-\frac{2c^2}{a^2+b^2+c^2}\biggr)^2=1.
\end{align*}
再由直线$L$的任意性知，满足条件的$Q$点到$M(0,0,1)$的距离恒为$1$。
\end{solution}

\begin{example}
已知方程$ f(x-2y,y-3z)=0 $。
\begin{enumerate}[(1)]
\item 利用柱面的定义，验证在通常情况下它表示空间中的一个柱面，并求出其母线方向。
\item 求出此柱面与母线方向垂直的一条准线。
\end{enumerate}
\end{example}
\begin{proof}
\begin{enumerate}[(1)]
\item 设向量$\overrightarrow{n}=\left( a,b,c \right)$，则将曲面$C:f\left( x-2y,y-3z \right) =0$沿$\overrightarrow{n}$的方向平移$t$个单位长得新曲面$C_t$方程为
\begin{align*}
f\left( \left( x+ta \right) -2\left( y+tb \right),\left( y+tb \right) -3\left( z+tc \right) \right) =f\left( \left( x-2y \right) +t\left( a-2b \right),\left( y-3z \right) +t\left( b-3c \right) \right) =0,\quad \forall t\in \mathbb{R}.
\end{align*}
取$c=1,b=3,a=6$，则此时$\overrightarrow{n}=\left( 6,3,1 \right)$，并且
\begin{align*}
C_t:f\left( \left( x+ta \right) -2\left( y+tb \right),\left( y+tb \right) -3\left( z+tc \right) \right) =f\left( x-2y,y-3z \right) =0,\quad \forall t\in \mathbb{R}.
\end{align*}
故曲面$C$和$C_t$重合，再由$t$的任意性知曲面$C$沿$\overrightarrow{n}=\left( 6,3,1 \right)$方向平移任意长得到的曲面还是$C$，故曲面$C$就是柱面，母线方向为$\overrightarrow{n}=\left( 6,3,1 \right)$。

\item 令平面$\pi :6x+3y+z=0$，则平面$\pi$与$\overrightarrow{n}=\left( 6,3,1 \right)$垂直。故曲线
\begin{align*}
L:\begin{cases}
6x+3y+z=0\\
f\left( x-2y,y-3z \right) =0
\end{cases}
\end{align*}
就是柱面与母线方向$\overrightarrow{n}=\left( 6,3,1 \right)$垂直的一条准线。
\end{enumerate}
\end{proof}

\begin{example}
曲线
\begin{align*}
\begin{cases}
xy+yz+zx=0,\\
x+y+z=3
\end{cases}
\end{align*}
在变换 $(x,y,z)\mapsto (x+y, y+2z, z+3x)$ 下变换为椭圆，求该椭圆的中心、法向量与面积。
\end{example}

\begin{solution}
设变换矩阵为
\begin{align}
A=\begin{pmatrix}
1&1&0\\
0&1&2\\
3&0&1
\end{pmatrix},
\quad \det A=7.\label{eq:detA}
\end{align}
原曲线满足
\begin{align}
xy+yz+zx=\frac{(x+y+z)^2-(x^2+y^2+z^2)}2,
\end{align}
代入 $x+y+z=3$ 得
\begin{align}
x^2+y^2+z^2=9.
\end{align}
因此原曲线是平面 $x+y+z=3$ 上的圆，圆心为 $c=(1,1,1)$，半径 $r=\sqrt{6}$，面积 $S_0=6\pi$.

变换后，椭圆中心为
\begin{align}
Ac=(2,3,4).
\end{align}
原平面法向量 $n=(1,1,1)$，像平面法向量 $m$ 满足 $A^T m=n$，解得
\begin{align}
m=\frac17(4,3,1),
\end{align}
故可取法向量为 $(4,3,1)$.

由面积放大公式（见下面命题），面积放大倍数为
\begin{align}
|\det A|\frac{\|A^{-T}n\|}{\|n\|}
=7\cdot \frac{\sqrt{26}/7}{\sqrt3}
=\sqrt{\frac{26}{3}}.
\end{align}
所以椭圆面积
\begin{align}
S=6\pi\sqrt{\frac{26}{3}}=2\pi\sqrt{78}.
\end{align}
综上，椭圆的中心为 $(2,3,4)$，法向量可取 $(4,3,1)$，面积为 $2\pi\sqrt{78}$.
\end{solution}

\begin{proposition}\label{proposition:线性变换放大空间平面面积}
设 $A$ 为三阶可逆实矩阵，$\pi$ 为平面 $n\cdot x=d$（其中 $n\in\mathbb R^3$，$n\neq 0$）。线性变换 $y=Ax$ 将 $\pi$ 上的任意曲线所围成的区域变为另一平面上的区域，则像平面的法向量为 $A^{-T}n$，且该变换将任意曲线所围成的面积放大为原来的
\begin{align*}
|\det A|\frac{\|A^{-T}n\|}{\|n\|}
\end{align*}
倍.
\end{proposition}
\begin{proof}
取像平面上的点 $y$，对应原平面上的点 $x=A^{-1}y$。代入原平面方程得
\begin{align*}
n\cdot (A^{-1}y)=d,
\end{align*}
即
\begin{align*}
(A^{-T}n)\cdot y=d.
\end{align*}
因此像平面的一个法向量为 $A^{-T}n$.

任取原平面上两个不共线的向量 $u,v$，它们张成的平行四边形面积元为 $\|u\times v\|$。变换后为 $Au,Av$，面积元为 $\|Au\times Av\|$。利用恒等式
\begin{align*}
(Au)\times (Av)=\det(A)\,A^{-T}(u\times v),
\end{align*}
取模得
\begin{align*}
\|Au\times Av\|=|\det A|\,\|A^{-T}(u\times v)\|.
\end{align*}
因为 $u\times v$ 垂直于原平面，故平行于 $n$，设 $u\times v=\lambda n$，则
\begin{align*}
\|u\times v\|=|\lambda|\,\|n\|.
\end{align*}
于是
\begin{align*}
\|Au\times Av\|
=|\det A|\,|\lambda|\,\|A^{-T}n\|
=|\det A|\frac{\|A^{-T}n\|}{\|n\|}\,\|u\times v\|.
\end{align*}
所以任意面积元都按同一倍数放大，因而由曲线围成的任意区域面积也按该倍数放大.
\end{proof}

\begin{example}
曲线
\begin{align*}
\begin{cases}
xy+yz+zx=0,\\
x+y+z=3
\end{cases}
\end{align*}
在变换 $(x,y,z)\mapsto (x+y, y+2z, z+3x)$ 下变换为椭圆，求该椭圆的中心、法向量与面积。
\end{example}
\begin{solution}
设变换矩阵为
\begin{align*}
A=\begin{pmatrix}
1&1&0\\
0&1&2\\
3&0&1
\end{pmatrix},
\quad \det A=7.\label{eq:detA_02}
\end{align*}
原曲线满足
\begin{align*}
xy+yz+zx=\frac{(x+y+z)^2-(x^2+y^2+z^2)}2,
\end{align*}
代入 $x+y+z=3$ 得
\begin{align*}
x^2+y^2+z^2=9.
\end{align*}
因此原曲线是平面 $x+y+z=3$ 上的圆，圆心为 $c=(1,1,1)$，半径 $r=\sqrt{6}$，面积 $S_0=6\pi$.

变换后，椭圆中心为
\begin{align*}
Ac=(2,3,4).
\end{align*}
由\refpro{proposition:线性变换放大空间平面面积}知,原平面法向量 $n=(1,1,1)$，像平面法向量 $m$ 满足 $A^T m=n$，解得
\begin{align*}
m=\frac17(4,3,1),
\end{align*}
故可取法向量为 $(4,3,1)$.

由\refpro{proposition:线性变换放大空间平面面积}，面积放大倍数为
\begin{align*}
|\det A|\frac{\|A^{-T}n\|}{\|n\|}
=7\cdot \frac{\sqrt{26}/7}{\sqrt3}
=\sqrt{\frac{26}{3}}.
\end{align*}
所以椭圆面积
\begin{align*}
S=6\pi\sqrt{\frac{26}{3}}=2\pi\sqrt{78}.
\end{align*}
综上，椭圆的中心为 $(2,3,4)$，法向量可取 $(4,3,1)$，面积为 $2\pi\sqrt{78}$.
\end{solution}


























































































































\end{document}